% Constructive section for "The Unchosen Adjacency" (revised)
% Assumes the definitions file. Discord timeout checked against Discord's support documentation.
% Meta comparators rest partly on secondary sources; see the bibliography notes.

\section{Separating Participation from Perception}\label{sec:moderation}

The applications section found that a block removes perception along with participation, that the controls found on personal Facebook profiles do not separate the two, and that other Meta surfaces do, silently in one case. This section describes what a transparent alternative looks like, in order to show that clean realization is a concrete design target and not an abstraction.

\subsection{Two kinds of relation}

Partition the relations that concern another person $v$ into two classes.

\begin{definition}[Participation and perception]
Let $R_u^{\mathrm{part}}(v)$ be the relations by which $v$ may act within $u$'s space: commenting, replying, tagging, messaging, and mentioning. Let $R_u^{\mathrm{perc}}(v)$ be the relations by which $v$ may observe $u$'s space: viewing posts, viewing responses, and viewing the norms in force.
\end{definition}

The partition is analytic. It says that the two classes are independently controllable in principle, and it identifies which relation a given action removes. A block removes both classes at once. The constructive proposal is an action that removes the participation the user intends to bar and preserves perception.

\begin{definition}[Asymmetric separation]
Let $P\subseteq R_u^{\mathrm{part}}(v)$ be the participation relations the user intends to bar. An action $a$ effects \emph{asymmetric separation of $v$ with respect to $P$} if
\[
E(a)\;\supseteq\;\{(r,-) : r\in P\}
\quad\text{and}\quad
E(a)\cap\{(r,-): r\in R_u^{\mathrm{perc}}(v)\}=\varnothing .
\]
\end{definition}

Defining separation relative to $P$ matters: the intention may be only to stop commenting, and an action that also removes messaging, tagging, and mentioning would bundle unwanted loss with the intended one.

Where a comparator exists, it supports feasibility. Discord's timeout, as documented by the service, leaves a member able to read while barring them from sending \cite{discord-timeout}. The comparator is a moderator power and not a user power, so it shows that the separation can be built and operated at scale and does not by itself show that a personal-profile equivalent has been deployed elsewhere.

\begin{remark}
Comparators on other Meta surfaces show that the separation is buildable within one company's products: temporary suspension of group members \cite{fb-groups}, bans on Pages that leave the Page visible \cite{page-ban}, and Instagram's Restrict, which hides a person's comments from others unless approved \cite{ig-restrict}. On personal Facebook profiles the controls found differ. The Restricted list limits what a person sees, namely public information and tagged posts \cite{fb-restricted}, and comment limits set an audience and not a person \cite{fb-comments}. Secondary guides that describe a Facebook Restrict hiding comments appear to conflate it with Instagram's feature. Instagram's Restrict is silent, so the person is neither told which norm was applied nor asked to revise, and no reason or duration is described. It is an indefinite reversible state and not a time-bounded suspension, and so it is a concealment tool more than a correction tool.
\end{remark}

\subsection{What the user values}

The case for retained perception does not require that seeing good conduct reforms the person who sees it. That is an empirical claim and may be false in many cases. The framework asks only that the user place value on the retained relation, that is, $V_u(r)\ge\lambda_{\tau(r)}$ for $r\in R_u^{\mathrm{perc}}(v)$. A user may value it for several reasons that do not depend on reform: keeping a shared space open, avoiding the appearance of hiding from a person, or preserving the possibility of later reconciliation. The reform hypothesis is one reason among these and is not a premise.

\subsection{A graduated moderation sequence}

Moderation in ordinary social life is graduated. A person is asked to stop, then asked to leave the conversation, then asked to leave the gathering, and only in the last case excluded from the community. The interface can be described by the same structure. An illustrative sequence for conduct in comment spaces has two parts.

\emph{Corrective stages}, which differ in how they treat an utterance and not in what access the person retains:
\begin{enumerate}
  \item \textbf{Pre-posting prompt.} Before submission, a notice states the norm in force and asks the author whether to revise. No relation is removed.
  \item \textbf{Return with explanation.} A posted comment is returned to its author with a stated reason and a request to remove or revise it. The comment is hidden pending response.
  \item \textbf{Removal with notice.} The comment is removed and a public note states the norm applied, without naming the author.
\end{enumerate}

\emph{Access restrictions}, which differ in how much participation or perception they remove:
\begin{enumerate}
  \setcounter{enumi}{3}
  \item \textbf{Temporary thread exclusion.} The person may not comment on the user's posts for a stated period, while retaining view.
  \item \textbf{Restricted interaction.} A longer or wider exclusion of participation, again preserving view.
  \item \textbf{Block.} All relations are removed, reserved for stalking, threats, persistent harassment, or cases where visibility itself creates danger.
\end{enumerate}

The corrective stages form a workflow rather than a chain, and the access restrictions form a ladder that can be ordered exactly.

\begin{definition}[Restriction set and ladder]
For an access restriction $a$ and target $v$, let
\[
Q_v(a)=\{\,r\in R_u^{\mathrm{part}}(v)\cup R_u^{\mathrm{perc}}(v) : (r,-)\in E(a)\,\}
\]
be its \emph{restriction set}, the participatory and perceptual relations it revokes. Order access restrictions by $a\preceq a'$ if $Q_v(a)\subseteq Q_v(a')$. A \emph{restriction ladder} is a chain $a_0\preceq a_1\preceq\cdots\preceq a_n$ in this order. Notices, reasons, and durations do not enter $Q_v$; they are intended changes in their own right.
\end{definition}

\begin{remark}
Complete effect sets are not ordered by inclusion. A prompt adds a warning, a removal deletes an utterance already made, a timeout disables future participation, and a block removes participation and perception, so their full signed effect sets are generally incomparable. Ordering by restriction sets isolates the dimension along which access severity increases. The corrective stages are not ordered by $Q_v$ because they restrict an utterance and not a person's access. Enlarging $Q_v$ to utterance-level relations and time-indexed relations would order them, at a cost in machinery that the argument does not need.
\end{remark}

\begin{definition}[Local collapse]
For an intended restriction $S\subseteq R_u^{\mathrm{part}}(v)\cup R_u^{\mathrm{perc}}(v)$, let $\mathcal{A}_S=\{a\in\mathcal{A} : S\subseteq Q_v(a)\}$ be the available access restrictions that cover it. If $\mathcal{A}_S=\varnothing$, the restriction is unrealizable at $\mathcal{A}$. Otherwise, for $a\in\mathcal{A}_S$ the \emph{excess restriction} is
\[
\Sigma(S,a)=Q_v(a)\setminus S .
\]
The restriction $S$ \emph{collapses} at $\mathcal{A}$ if $\mathcal{A}_S\neq\varnothing$ and $\Sigma(S,a)\neq\varnothing$ for every $a\in\mathcal{A}_S$. The intended restriction is then forced upward, and $\Sigma$ measures how far.
\end{definition}

\begin{remark}
Collapse is local to an intended restriction. An interface may offer several finer controls and still collapse for a particular intention, because no offered action covers that intention without excess. This avoids the false claim that an interface offers only the top rung. Collapse is the overreach failure of the access restrictions, and the omission of notice, reason, or duration is the underrealization failure. When $S\subseteq R_u^{\mathrm{part}}(v)$ collapses and every covering action has excess containing a perceptual relation $r$ with $V_u(r)\ge\lambda_{\tau(r)}$, the excess is consequential loss, and if a covering action with no consequential excess is feasible at $\mathcal{A}^{*}$, the intention is withheld in the earlier sense.
\end{remark}

\subsection{Norm-stating removals}

The present arrangement also privatizes the lesson. When a user blocks someone, other participants see only that the disruption ended. They do not learn which norm was applied, and the person moderated cannot observe how the conversation proceeds. A removal notice of the form ``a comment was removed because it prescribed another participant's private conduct instead of addressing the topic'' would make the norm public without exposing the individual. Moderation then also articulates a rule, which is a relation the community may want and cannot obtain from a silent block. This is an effect of a third kind, neither association nor dissociation: the production of public institutional memory. It can enter the intended set as an added relation, and an interface whose only covering action is a silent block cannot realize it.

\subsection{Limits}

Three limits should be stated. The classification of a comment as addressing character rather than argument is contestable and error-prone, so pre-posting prompts risk false positives and paternalism; the proposal is a design possibility and not a claim about the accuracy of any classifier. Graduated systems have costs in complexity and in the labor of the person applying them, which returns to the simplicity objection above. And asymmetric separation does not prevent circumvention, since a restricted person may route around it through direct messages or another account, though neither does a block. The section does not claim that the sequence is optimal. It claims that the lower stages are realizable, that their absence is a design fact, and that their absence is what turns each mild intention into an exile.
